\documentclass[11pt,a4paper]{article}
\usepackage[T1]{fontenc}
\usepackage[utf8]{inputenc}
\usepackage{lmodern}
\usepackage[czech]{babel}
\usepackage[margin=1.7cm]{geometry}
\usepackage{amsmath,amssymb,mathtools}
\usepackage{xcolor,cancel,enumitem}
\usepackage{tikz,pgfplots}
\pgfplotsset{compat=1.18}
\usetikzlibrary{arrows.meta,calc,patterns}
\pagestyle{empty}
\setlength{\parindent}{0pt}
\setlength{\parskip}{4pt}
\newcommand{\vysl}[1]{\boxed{\color{red!75!black}#1}}
\newcommand{\pozn}[1]{{\color{red!75!black}#1}}
\newcommand{\dx}{\,\mathrm{d}x}
\newcommand{\dt}{\,\mathrm{d}t}
\newcommand{\subst}[1]{\left|\begin{array}{l}#1\end{array}\right|}
\DeclareMathOperator{\arctg}{arctg}
\DeclareMathOperator{\arccotg}{arccotg}
\DeclareMathOperator{\tg}{tg}
\DeclareMathOperator{\cotg}{cotg}
\usepackage{tcolorbox}
\newtcolorbox{pomocny}{colframe=red!75!black, colback=white, coltitle=white, fonttitle=\bfseries, title={Pomocný výpočet:}, sharp corners}
\newcommand{\dy}{\,\mathrm{d}y}
% znaménkové schéma: \znam{bod/popis/typ,...}{znaménka v intervalech}
% typ: o = vyloučený bod (kroužek), x = obyčejný bod
\newcommand{\znam}[2]{%
\begin{tikzpicture}[baseline=-3pt]
\def\n{0}
\foreach \p/\l/\t [count=\i] in {#1} {\xdef\n{\i}}
\pgfmathsetmacro{\L}{1.6*(\n+1)}
\draw (0,0) -- (\L,0);
\foreach \p/\l/\t [count=\i] in {#1} {
  \pgfmathsetmacro{\xx}{1.6*\i}
  \ifnum\pdfstrcmp{\t}{o}=0 \draw[fill=white] (\xx,0) circle (2.5pt); \draw (\xx-0.09,0) -- (\xx+0.09,0) (\xx,-0.09) -- (\xx,0.09);
  \else \draw (\xx,0.1) -- (\xx,-0.1); \fi
  \node[below=2pt] at (\xx,0) {\small $\l$};
}
\foreach \s [count=\i] in {#2} {
  \pgfmathsetmacro{\xx}{1.6*\i-0.8}
  \node[above] at (\xx,0) {$\s$};
}
\end{tikzpicture}}

\newcommand{\krok}[1]{{\small\color{blue!60!black}#1}}
\newcommand{\kroku}[1]{\text{\small\color{blue!60!black}#1}}

\begin{document}

\section*{Determinanty -- řešení}

\textbf{a)}
\[
\begin{vmatrix} -3 & 9 & 3 & 6\\ -5 & 8 & 2 & 7\\ 4 & -5 & -3 & -2\\ 7 & -8 & -4 & -5 \end{vmatrix}
= \underbrace{3}_{\substack{\kroku{vytknutí 3 z}\\ \kroku{prvního řádku}}}\cdot
\begin{vmatrix} -1 & 3 & 1 & 2\\ -5 & 8 & 2 & 7\\ 4 & -5 & -3 & -2\\ 7 & -8 & -4 & -5 \end{vmatrix} =
\]
\krok{Úprava 1. sloupce: k 2. řádku přičteme $(-5)$-násobek 1. řádku, ke 3. řádku $4$-násobek a ke 4. řádku $7$-násobek 1. řádku.}
\[
= 3\cdot\begin{vmatrix} -1 & 3 & 1 & 2\\ 0 & -7 & -3 & -3\\ 0 & 7 & 1 & 6\\ 0 & 13 & 3 & 9 \end{vmatrix} =
\]
\[
= \underbrace{3\cdot(-1)\cdot(-1)^{1+1}}_{-3}\cdot
\begin{vmatrix} -7 & -3 & -3\\ 7 & 1 & 6\\ 13 & 3 & 9 \end{vmatrix} =
\]
\krok{Rozvoj podle 1. sloupce; pak úprava 2. sloupce: k 1. řádku přičteme $3$-násobek 2. řádku, ke 3. řádku $(-3)$-násobek 2. řádku.}
\[
= -3\cdot\begin{vmatrix} 14 & 0 & 15\\ 7 & 1 & 6\\ -8 & 0 & -9 \end{vmatrix}
= \underbrace{-3\cdot1\cdot(-1)^{2+2}}_{-3}\cdot\begin{vmatrix} 14 & 15\\ -8 & -9 \end{vmatrix}
= -3\big(14\cdot(-9)-15\cdot(-8)\big) =
\]
\krok{Rozvoj podle 2. sloupce.}
\[
= -3\,(-126+120) = -3\cdot(-6) = \vysl{18}
\]

\medskip
\hrule
\medskip

\textbf{b)}
\[
\begin{vmatrix} -5 & 0 & -1 & -1\\ -1 & 3 & 0 & 2\\ -1 & 1 & 2 & 0\\ 2 & 1 & 1 & 1 \end{vmatrix}
= \begin{vmatrix} -3 & 1 & 0 & 0\\ -1 & 3 & 0 & 2\\ -5 & -1 & 0 & -2\\ 2 & 1 & 1 & 1 \end{vmatrix} =
\]
\krok{Úprava 3. sloupce: k 1. řádku přičteme 4. řádek, ke 3. řádku $(-2)$-násobek 4. řádku. Pak rozvoj podle 3. sloupce.}
\[
= \underbrace{1\cdot(-1)^{4+3}}_{-1}\cdot\begin{vmatrix} -3 & 1 & 0\\ -1 & 3 & 2\\ -5 & -1 & -2 \end{vmatrix}
= -1\cdot\begin{vmatrix} -3 & 1 & 0\\ -1 & 3 & 2\\ -6 & 2 & 0 \end{vmatrix} =
\]
\krok{Úprava 3. sloupce: ke 3. řádku přičteme 2. řádek. Pak rozvoj podle 3. sloupce.}
\[
= \underbrace{-1\cdot2\cdot(-1)^{2+3}}_{2}\cdot\begin{vmatrix} -3 & 1\\ -6 & 2 \end{vmatrix}
= 2\big(\underbrace{-3\cdot2-1\cdot(-6)}_{0}\big) = \vysl{0}
\]

\end{document}
